\section{Lemat o punkcie stałym (diagonalizacja)}\label{sec:diagonalizacja}
\begin{lemat}[Lemat diagonalny / o punkcie stałym]\label{res:8:1}
Niech $\psi(x)$ będzie formułą, której jedyną dopuszczalną zmienną wolną jest $x$. Istnieje zdanie $G$ takie, że $Q\vdash G\leftrightarrow\psi(\ulcorner G\urcorner)$. Ta sama równoważność jest dowodliwa w każdej $T\supseteq Q$, niezależnie od efektywności jej aksjomatów. Zdanie $G$ „mówi o sobie”, że ma własność $\psi$; nie jest to dosłowna równość napisów, lecz dowodliwa równoważność.
\end{lemat}
\begin{proof}
Ustalamy wyróżnioną zmienną $x$ i definiujemy funkcję diagonalizacji
\[
d(n):=\ulcorner\theta(\overline n)\urcorner,
\]
gdy $n$ koduje formułę $\theta(x)$ bez innych zmiennych wolnych; na pozostałych wejściach kładziemy $d(n)=0$. Jest to funkcja p.r. na mocy Lematu~\ref{res:6:5}(c): podstawiamy numerał $\overline n$ w formule o kodzie $n$. Twierdzenie~\ref{res:7:6} daje formułę $\varphi_d(x,y)$ z
\[
Q\vdash\forall y\,(\varphi_d(\overline n,y)\leftrightarrow y=\overline{d(n)})
\]
dla każdego standardowego $n$.

Przed podstawieniami przemianowujemy zmienne związane, aby nowa zmienna $y$ nie została przechwycona. Definiujemy $\delta(x):=\exists y\,(\varphi_d(x,y)\land\psi(y))$, $m:=\ulcorner\delta(x)\urcorner$ i $G:=\delta(\overline m)$. $G$ jest zdaniem, ponieważ wszystkie wolne wystąpienia $x$ zastąpiliśmy domkniętym numerałem. Z definicji $d(m)=\ulcorner\delta(\overline m)\urcorner=\ulcorner G\urcorner$, więc reprezentowalność daje
\[
Q\vdash\forall y\,(\varphi_d(\overline m,y)\leftrightarrow y=\ulcorner G\urcorner).
\]
Podstawiając do definicji $G$, uzyskujemy $Q\vdash G\leftrightarrow\exists y\,(y=\ulcorner G\urcorner\land\psi(y))\leftrightarrow\psi(\ulcorner G\urcorner)$. To jest szukana równoważność. Nie użyto żadnej hipotezy o aksjomatach $T$.
\end{proof}

